\unitopen{09}
\begin{candos}
\item kupić podstawowe rzeczy w sklepie, podając ilość i rozmiar;
\item odmienić czasowniki z zamianą samogłoski w drugiej i trzeciej osobie (\de{nehmen, geben, fahren, lesen, tragen});
\item zapłacić gotówką i kartą i zrozumieć kwotę;
\item opisać ubranie kolorem i rozmiarem.
\end{candos}

\section{Słownictwo}
\vocabtable{09}

\section{Czasowniki z zamianą samogłoski}
W drugiej i trzeciej osobie liczby pojedynczej niektóre czasowniki zmieniają samogłoskę rdzenia. Pozostałe osoby
(\emph{ich}, \emph{wir}, \emph{ihr}, \emph{sie}) są regularne.

\begin{center}\small
\begin{tabular}{@{}lllll@{}}
\toprule
\textsf{\textbf{Zmiana}} & \textsf{\textbf{Bezokolicznik}} & \textsf{\textbf{ich}} & \textsf{\textbf{du}} & \textsf{\textbf{er/sie/es}}\\
\midrule
\de{e $\rightarrow$ i} & \de{nehmen} & \de{nehme} & \de{nimmst} & \de{nimmt}\\
& \de{geben} & \de{gebe} & \de{gibst} & \de{gibt}\\
& \de{sprechen}, \de{essen} & \de{spreche}, \de{esse} & \de{sprichst}, \de{isst} & \de{spricht}, \de{isst}\\
\de{e $\rightarrow$ ie} & \de{lesen}, \de{sehen} & \de{lese}, \de{sehe} & \de{liest}, \de{siehst} & \de{liest}, \de{sieht}\\
\de{a $\rightarrow$ ä} & \de{fahren} & \de{fahre} & \de{fährst} & \de{fährt}\\
& \de{tragen} & \de{trage} & \de{trägst} & \de{trägt}\\
\bottomrule
\end{tabular}
\end{center}
Nie zgadujesz, który czasownik się zmienia: uczysz się \emph{bezokolicznika razem z formą \de{du}}. Fiszki w \code{anki/}
mają formy w kolumnie uwag. \de{Es gibt} + biernik znaczy „jest/są”: \de{Es gibt einen Supermarkt in der Nähe.}

\section{Ilości, ceny, kolory}
Po miarze nie ma przyimka, rzeczownik stoi bez końcówki mnogiej po rzeczach policzalnych w mianowniku: \de{ein Kilo Äpfel},
\de{zwei Liter Milch}, \de{drei Packungen Nudeln}, \de{zweihundert Gramm Käse}. \de{Wie viel?} dla niepoliczalnych
(\de{Wie viel Käse?}), \de{Wie viele?} dla policzalnych (\de{Wie viele Äpfel?}). Kolory: \de{rot, blau, grün, schwarz, weiß, gelb}.
Przymiotnik w roli orzecznika nie ma końcówki: \de{Die Jacke ist rot.}

\begin{falle}
Rozmiar ubrań w DACH: \de{Größe 38}, \de{Größe M}, a buty \de{Schuhgröße 42}. Nie mylisz ich z europejskim
(\de{M} = \de{medium}). I pułapka językowa: \de{Ich fahre zur Arbeit} (jadę) vs \de{Ich gehe zu Fuß} (idę): niemiecki
rozróżnia \emph{jechać} (\de{fahren}) i \emph{iść} (\de{gehen}), podobnie jak polski, ale \de{fahren} dotyczy każdego pojazdu,
także roweru: \de{Ich fahre mit dem Fahrrad}.
\end{falle}

\begin{szwajcaria}
Płatność kartą jest w Szwajcarii powszechna, ale małe sklepy lub targi mogą brać tylko gotówkę;
szwajcarskie kasy zaokrąglają do 5 Rappen. Torebka w sklepie bywa płatna (\de{Brauchen Sie eine Tüte?} $\rightarrow$ w CH
częściej \de{Säckli}). Znani sprzedawcy: \de{Migros}, \de{Coop}; zakupy niedzielne są mocno ograniczone
\verify{(godziny otwarcia zależą od kantonu)}.
\end{szwajcaria}

\begin{praca}
Zamówienie sprzętu do biura: \de{Wir brauchen zehn Laptops und fünf Bildschirme.} \de{Was kostet ein Bildschirm
mit Lieferung?} \de{Wann können Sie liefern?} Po niemiecku rozmiar \de{Größe M} dotyczy ubrania, a rozdzielczość to \de{die Auflösung}.
\end{praca}

\section{Dialogi}
\textbf{Dialog 1 --- w sklepie z ubraniami}
\begin{dialog}
\turn{A}{Guten Tag. Kann ich Ihnen helfen?}{Dzień dobry. Czy mogę w czymś pomóc?}
\turn{B}{Ja, ich suche einen Pullover. Haben Sie den auch in Größe M?}{Tak, szukam swetra. Czy macie go w rozmiarze M?}
\turn{A}{Ja, den gibt es in Blau und Grün. Welche Farbe möchten Sie?}{Tak, jest w niebieskim i zielonym. Jaki kolor Pan chce?}
\turn{B}{Ich nehme den blauen. Was kostet er?}{Wezmę niebieski. Ile kosztuje?}
\turn{A}{Fünfzig Euro. Bar oder mit Karte?}{Pięćdziesiąt euro. Gotówką czy kartą?}
\turn{B}{Mit Karte, bitte. Das ist alles.}{Kartą, proszę. To wszystko.}
\end{dialog}

\section{Ćwiczenia}
\begin{cwiczenia}
\exitem{Odmień \de{nehmen} dla \de{du}.}{\de{du nimmst}}
\exitem{Odmień \de{fahren} dla \de{er}.}{\de{er fährt}}
\exitem{Odmień \de{lesen} dla \de{du}.}{\de{du liest}}
\exitem{Odmień \de{geben} dla \de{er}.}{\de{er gibt}}
\exitem{Odmień \de{tragen} dla \de{du}.}{\de{du trägst}}
\exitem{Uzupełnij: \de{Wir \dots{} den Zug.} (nehmen)}{\de{nehmen} (w liczbie mnogiej bez zmiany samogłoski)}
\exitem{Uzupełnij: \de{Ich möchte ein \dots{} Äpfel.} (kilogram)}{\de{Kilo}}
\exitem{Przeczytaj: 200 g Käse.}{\de{zweihundert Gramm Käse}}
\exitem{Wybierz \de{wie viel} czy \de{wie viele}: \dots{} Brötchen?}{\de{Wie viele Brötchen?}}
\exitem{Wybierz \de{wie viel} czy \de{wie viele}: \dots{} Käse?}{\de{Wie viel Käse?}}
\exitem{Zapytaj, czy mają rozmiar M.}{\de{Haben Sie das in Größe M?}}
\exitem{Powiedz, że coś jest za drogie.}{\de{Das ist zu teuer.}}
\exitem{Przetłumacz: „Zapłacę kartą.”}{\de{Ich zahle mit Karte.}}
\end{cwiczenia}

\weekend{Zrób w domu \emph{listę zakupów} po niemiecku (\de{ein Kilo Äpfel, zwei Liter Milch, \dots}) z ilościami. Zasymuluj
z AI rozmowę w sklepie: poproś o trzy rzeczy, zapytaj o cenę, odrzuć jedną jako zbyt drogą i zapłać.}
